%%%%%%%%%% source scan idx 13 | folio 6 %%%%%%%%%%

1.10. \underline{本节余下部分只考虑例 1.1 中的情形。}

Lichtenbaum [5] 已证明，在此情形下 $X$ 自动在 $S$ 上射影，因而定义复形 $Rf^{!}\underline{O}_S$ 并无困难（[4] III 8.7，第~190 页）。由于 $f$ 的相对维数为 1，且 $X$ 与 $S$ 均正则（Gorenstein 即已足够），存在 $X$ 上的可逆层 $\omega_{X/S}$，在唯一同构意义下唯一（[4] V 8.3 及 V 9.1），使得

\begin{equation}\tag{1.10.1}
Rf^{!}\underline{O}_S \;\simeq\; \omega_{X/S}[1] \quad .
\end{equation}
用 $K$ 表示满足 $\omega_{X/S}\simeq\underline{O}(K)$ 的任一除子（有效与否均可）。

设 $D$ 是支集包含于 $Y$ 的除子 $\geq0$，$N$ 是其余法层：

\begin{center}
\begin{tikzcd}[column sep=large, row sep=large]
D \arrow[rr, "i"] \arrow[dr, swap, "g"] & {} & X \arrow[dl, "f"] \\
{} & S & {}
\end{tikzcd}
\end{center}
有

\begin{equation}\tag{1.10.2}
Rg^{!}\underline{O}_S
=Ri^{!}Rf^{!}\underline{O}_S
=Ri^{!}\omega_{X/S}[1]
=\omega_{X/S}\big|_D\otimes N^{\vee}[0].
\end{equation}
由对偶定理（[4] V 11.1，第~210 页），

\begin{equation}\tag{1.10.3}
Rg_*Rg^{!}\underline{O}_S
\;\simeq\;
\operatorname{RHom}(Rg_*\underline{O}_D,\underline{O}_S) \quad .
\end{equation}
此外，对每个支集包含于闭点的 $\underline{O}_S$-模 $M$，

\begin{equation}\tag{1.10.4}
\chi\bigl(\operatorname{RHom}(M,\underline{O}_S)\bigr)=-\chi(M),
\end{equation}
所以这里利用（1.4.4）和（1.10.3）得到

\begin{equation*}
\chi(Rg^{!}\underline{O}_S)=-\chi(\underline{O}_D) \quad .
\end{equation*}
又由（1.4.6），

\begin{equation*}
\chi(Rg^{!}\underline{O}_S)
=\chi(\underline{O}_D)+\deg_D(Rg^{!}\underline{O}_S),
\end{equation*}

%%%%%%%%%% source scan idx 14 | folio 7 %%%%%%%%%%

结合（1.10.2）便得

\begin{equation*}
-2\chi(\underline{O}_D)
=\deg_D(Rg^{!}\underline{O}_S)
=\deg_D(\omega_{X/S}|_D\otimes N^{\vee})
=\deg_D(\omega_{X/S})-\deg_D(N),
\end{equation*}
由 1.6(ii) 和 (v)，可改写为如下命题。

\underline{命题} 1.11. \underline{在（1.1）的假设下，对支集包含于 $Y$ 的任一除子 $D\geq0$，有}

\begin{equation*}
\chi(\underline{O}_D)=-\frac12(D,D+K) \quad .
\end{equation*}

1.12. 用 $d$ 表示 $X_s=f^{-1}(s)$ 各不可约分支的几何重数的最大公因数；也就是说，是 $f$ 的几何特殊纤维各泛点处局部环长度的最大公因数。若 $X_\eta$ 有有理点，等价地，若 $f:X\longrightarrow S$ 有截面 $x$，则 $d=1$；更确切地说，$f$ 在 $x(s)$ 处光滑。事实上，$x$ 是正则概形之间的闭浸入，因而是正则浸入，可应用 EGA IV 17.12.1。

下面给出 Raynaud [9] 的一个定理的特殊情形的 Artin 证明；Raynaud 的定理只假设 $d$ 与 $k$ 的特征 $p$ 互素。

\underline{定理} 1.13. \underline{沿用上述记号，若 $Y$ 几何连通且 $d=1$，则 $X$ 在 $S$ 上为 $0$ 维上同调平坦的}（EGA III 7.8.1）；\underline{更确切地说}（EGA III 7.6.9(ii)），\underline{有}

\begin{equation*}
k\xrightarrow{\;\sim\;}H^0(X_s,\underline{O}_{X_s}).
\end{equation*}
容易归结到 $S$ 严格局部的情形。由于 $Y$ 几何连通，$H^0(Y,\underline{O}_Y)$ 是 $k$ 的纯不可分扩张；又因为 $(d,p)=1$，$Y$ 至少有一个不可约分支是

%%%%%%%%%% source scan idx 15 | folio 8 %%%%%%%%%%

几何约化的。于是容易得到 $H^0(Y,\underline{O}_Y)=k$。

设 $D$ 是 $X$ 上的除子，满足

\begin{equation*}
Y\leq D<X_s \quad .
\end{equation*}
由于 $d=1$，$D$ 不是 $X_s$ 的有理倍数；故由 1.8，

\begin{equation*}
(D,X_s-D)=-(D,D)>0 \quad .
\end{equation*}
所以除子 $Y$ 有一个不可约分支 $E$，使得

\begin{equation*}
0<E\leq X_s-D \quad\text{且}\quad (D,E)>0 \quad .
\end{equation*}
考虑图表

\begin{center}
\begin{tikzcd}[column sep=large, row sep=large]
0 \arrow[r] & \underline{O}(-D-E) \arrow[r] \arrow[d, "\varphi"] &
\underline{O} \arrow[r] \arrow[d, equal] &
\underline{O}_{D+E} \arrow[r] \arrow[d, "\psi"] & 0 \\
0 \arrow[r] & \underline{O}(-D) \arrow[r] &
\underline{O} \arrow[r] & \underline{O}_D \arrow[r] & 0
\end{tikzcd}
\end{center}
蛇引理给出 $\operatorname{Ker}(\psi)$ 与 $\operatorname{Coker}(\varphi)$ 的同构。这个余核恰为 $\underline{O}(-D)$ 在 $E$ 上的限制，故有正合列

\begin{equation*}
0\longrightarrow\underline{O}(-D)|_E
\longrightarrow\underline{O}_{D+E}
\longrightarrow\underline{O}_D\longrightarrow0 \quad .
\end{equation*}
由构造，$\deg_E(\underline{O}(-D)|_E)=-(D,E)<0$，所以 $H^0(\underline{O}(-D)|_E)=0$，且限制映射

\begin{equation*}
0\longrightarrow H^0(\underline{O}_{D+E})
\longrightarrow H^0(\underline{O}_D)
\end{equation*}
为单射。于是可用归纳法构造一列除子

\begin{equation*}
Y=D_0\leq D_1\leq\cdots\leq D_n=X_s
\end{equation*}
使得 $H^0(\underline{O}_{D_i})$ 单射到 $H^0(\underline{O}_{D_{i-1}})$。因此 $H^0(\underline{O}_{X_s})=H^0(\underline{O}_Y)=k$，1.13 得证。

%%%%%%%%%% source scan idx 16 | folio 9 %%%%%%%%%%

下述事实在 [5] 或 [10] 中有证明。

\underline{定理} 1.14（Castelnuovo、Enriques、\ldots）。\underline{在（1.1）的假设下，设 $D$ 是包含于 $Y$ 的约化不可约除子。假设 $H^1(D,\underline{O}_D)=0$，且置 $k'=H^0(D,\underline{O}_D)$ 后有}

\begin{equation*}
(D,D)=-[k':k] \quad .
\end{equation*}
\underline{则 $D$ 是 $k'$ 上的射影直线，并且存在正则 $S$-概形 $X_1$、$X_1$ 上剩余域为 $k'$ 的闭点 $x_1$，以及从 $X$ 到将 $X_1$ 在 $x_1$ 处爆破所得的 $S$-概形的同构，使 $D$ 对应于 $x_1$ 上的纤维。}

具有这些性质的除子 $D$ 称为\underline{第一类例外曲线（exceptional curve of the first kind）}。

我忍不住要证明下面这个 Néron 的命题（[8]，命题 1$'$，第~90 页）。

\underline{命题} 1.15. \underline{假设 $X_\eta$ 是椭圆曲线（特别地，它有截面），且 $X$ 不含第一类例外曲线。则沿用（1.10.1）的记号，典范映射}

\begin{equation*}
f^*f_*\omega_{X/S}\xrightarrow{\;\sim\;}\omega_{X/S}
\end{equation*}
\underline{普遍为同构}。

由于 $d=1$（1.12），由（1.13）只须证明 $\omega_{X/S}$ 同构于 $\underline{O}_X$。因为 $\omega_{X/S}$ 在椭圆曲线 $X_\eta$ 上平凡，先验地知道

\begin{equation*}
\omega_{X/S}\simeq\underline{O}\Bigl(\sum_i n_iY_i\Bigr),
\end{equation*}

%%%%%%%%%% source scan idx 17 | folio 10 %%%%%%%%%%

其中 $Y_i$ 是 $Y$ 的各不可约分支；又因 $\operatorname{div}(\pi)\sim0$，可假设 $n_i$ 为正。若 $\sum_i n_iY_i$ 与 $\operatorname{div}(\pi)$ 成比例，它就是 $\operatorname{div}(\pi)$ 的整数倍，证明即告完成。否则，

\begin{equation*}
\Bigl(\sum_i n_iY_i\Bigr)^2<0,
\end{equation*}
存在某个 $i$，使得

\begin{equation*}
\Bigl(Y_i,\sum_j n_jY_j\Bigr)<0,
\end{equation*}
便可应用下述引理。

\underline{引理} 1.16. \underline{在 1.1 的假设下，假定 $Y$ 连通，并且 $Y$ 可约或 $\chi(\underline{O}_Y)\leq0$。则 $Y$ 的一个不可约分支 $Y_i$ 是第一类例外曲线，当且仅当 $(Y_i,K)<0$。}

置 $k'=H^0(Y_i,\underline{O}_{Y_i})$。Riemann--Roch 公式 1.11 给出

\begin{equation*}
\chi(\underline{O}_{Y_i})
=[k':k]\bigl(1-\dim_{k'}H^1(Y_i,\underline{O}_{Y_i})\bigr)
=-\frac12\bigl((Y_i,Y_i)+(Y_i,K)\bigr) \quad .
\end{equation*}
若 $Y_i$ 是第一类例外曲线，即得 $(Y_i,K)=-[k':k]<0$。反之，若 $(Y_i,K)<0$，则 $\chi(\underline{O}_{Y_i})>0$，从而 $\chi(\underline{O}_{Y_i})=[k':k]$ 且 $Y$ 可约。故由 1.8，$(Y_i,Y_i)<0$。负整数 $(Y_i,Y_i)$ 和 $(Y_i,K)$ 均可被 $[k':k]$ 整除，所以必有

\begin{equation*}
(Y_i,Y_i)=(Y_i,K)=-[k':k] \quad .
\end{equation*}

\underline{习题}。(a) 在（1.1）的假设下，假定 $X_\eta$ 是某个椭圆曲线下的主齐性空间，且 $X$ 不含第一类例外曲线。证明，在 $X$ 上 Riemann--Roch 公式（1.11）具有如下形式

\begin{equation*}
\chi(\underline{O}_D)=-\frac12(D,D).
\end{equation*}

%%%%%%%%%% source scan idx 18 | folio 11 %%%%%%%%%%

(b) 再假定剩余域 $k(s)$ 代数闭。证明，$X_s$ 的不可约分支图是 Néron [8] 第~124--125 页第四栏所列某个图（情形 a 至 c8）的倍数。该图确定各约化不可约分支的亏格、它们在 $X_s$ 中的重数，以及 $(X_s)_{\mathrm{red}}$ 各奇点的类型和分布。

先处理 $X_s$ 不可约的情形。排除这种情形后，证明每个约化不可约分支都是自交数为 $-2$ 的射影直线。若 $X_s=\sum_i n_iD_i$ 是 $X_s$ 分解为约化不可约除子的表达式，将这个条件改写为

\begin{equation}\tag{*}
n_i=\frac12\sum_{j\ne i}n_j(D_i,D_j) \quad .
\end{equation}
另一方面，$X_s$ 连通，所以：

(**) 对任意 $i$ 和 $j$，存在链 $i=i_1,i_2,\ldots,i_n=j$，使 $(D_{i_k},D_{i_{k+1}})\ne0$。

现在求出 $(*)+(**)$ 的所有整数解。先证明，除了 Néron 的情形 b2 和 c2 外，相交数 $(D_i,D_j)$ 均为零或一。

